\chapter{Longitudinal and survival software as a decision model}
\label{exec17:s5}

A bank comparing repeated servicing policies needs more than a curve with a
confidence ribbon. An account can default, repay early, cease to be observed,
or receive a different action after its financial condition changes. A policy
that reduces default incidence can also shorten the profitable relationship.
A policy that improves every baseline risk stratum can have a larger late
hazard among its remaining accounts. Two programs can return numerically
similar objects named survival while counting different accounts. These
differences affect both the scientific question and the financial decision.

The longitudinal identification and sequential-estimation results developed
earlier supply the foundation. Here the question is how to preserve that
foundation when choosing and combining software. The argument proceeds from
the event clock and supported horizon to survivor selection, competing-event
simulation, the meaning of a modified action, and uncertainty in the resulting
monetary contrast. Its finite laws and monetary amounts are stipulated,
reproducible constructions. They provide exact comparison references and
expose otherwise hidden distinctions; they are not bank observations or new
executions of the external estimators.~\citep{exec17_s5_math}

\section{The decision clock and the unsupported tail}
\label{exec17:s5-horizon}

Let $j=1,\ldots,K$ index servicing periods. At the opening of a period the
bank observes the permitted history $H_j$, chooses an action $A_j$, and then
observes transitions and exits during that period. An intervention must specify
what the bank would do using the history available at that opening. A balance
posted after the decision, a later repayment classification, and a corrected
record unavailable to the decision maker cannot enter that history merely
because all three eventually appear in a warehouse. The data schema must
represent availability as well as the time to which a value refers.

Start with a single first-exit time $T$ taking values in the positive integers,
possibly with mass beyond the horizon. Under a specified policy $d$, define
$S^d(j)=P(T^d>j)$ and $S^d(0)=1$. Thus $S^d(j-1)$ is the probability of an
active account at the opening of period $j$. With payment for each period
entered, the expected restricted active time is
\begin{equation}
 R^d(K)=E\{\min(T^d,K)\}
       =\sum_{j=0}^{K-1}S^d(j).
 \label{exec17:s5-restricted-time}
\end{equation}
The identity follows account by account from
$\min(T,K)=\sum_{j=0}^{K-1}1\{T>j\}$ before taking expectations. This convention
earns a period's margin even if exit occurs at its end. A different payment
clock needs its own reward function. In particular, the same numerical area
must not be priced as both opening-period activity and completed-period
activity.

Suppose the observation and identification assumptions support the policy
survival probabilities only through period $m<K$. This may be an administrative
follow-up limit, rather than random censoring that an observation model could
remove. A fitted survival function beyond $m$ then supplies an extrapolation
assumption. Increasing simulation draws does not add observations at those
times. Without a tail restriction, monotonicity of survival gives the sharp
discrete bounds
\begin{equation}
 R^d(m)+S^d(m)\ \leq\ R^d(K)\ \leq\
 R^d(m)+(K-m)S^d(m).
 \label{exec17:s5-tail-bounds}
\end{equation}
To prove the lower bound, expand the remaining terms as
$S^d(m)+S^d(m+1)+\cdots+S^d(K-1)$ and use nonnegativity. To prove the upper
bound, bound every one of these $K-m$ terms by $S^d(m)$. The known term
$S^d(m)$ cannot be discarded. The continuous-time argument that allows all
remaining activity to disappear immediately after an observed time would
miss this whole discrete period.

Both endpoints are attainable. Keep the observed distribution through $m$
unchanged. For the lower endpoint, make all accounts still active at $m$ exit
at the end of period $m+1$. For the upper endpoint, keep them active through
the entire requested horizon. Mixing these two continuation laws preserves
every observed survival probability and fills the interval. The result is an
identified set conditional on the supported earlier survival law, not a
confidence interval for a point target. Sampling uncertainty about that earlier
law is a further problem.

For an exact six-period example, take the observed policy survival sequence
at times $0,1,2,3$ to be $(1,0.8,0.65,0.5)$ and the reference sequence to be
$(1,0.9,0.8,0.6)$. Then
\begin{align}
 R^d(6)&\in[2.95,3.95],& R^0(6)&\in[3.30,4.50],\notag\\
 R^d(6)-R^0(6)&\in[-1.55,0.65].
 \label{exec17:s5-tail-example}
\end{align}
These contrast endpoints are sharp when there is no additional cross-policy
restriction: select the lower continuation for one policy and the upper for
the other. Such potential event-time marginals can coexist, for example by
constructing each from a common uniform variable and its inverse distribution
function. A scientifically justified individual monotonicity restriction
would be additional information and could narrow the set.

At USD~20 margin per period entered and USD~3 upfront policy cost, the
incremental retention-margin value is therefore in $[-34,10]$ USD per eligible
account. This example prices activity alone; it does not silently include a
default loss or a repayment penalty. Those would require the corresponding
joint event probabilities. A decision to roll out the policy cannot be
justified by selecting the favorable end of this set. A bank with a more
restrictive accepted tail model could make a different, explicitly conditional
decision.~\citep{exec17_s5_math}

There are several substantive remedies. Shortening the requested horizon to
the supported period changes the business question but removes this tail
ambiguity. Following a public or generated cohort longer can demonstrate how
a proposed analysis would use additional follow-up, without supplying the
bank's missing future. An externally justified survival envelope or a specified
parametric tail gives conditional bounds or a point extrapolation. The
report should show how the monetary conclusion changes across those choices.
An automatic default horizon chosen from the largest event in the current
sample answers a fourth question, whose end point varies with the sample.
That can be useful internally, but it does not establish comparability between
policies, software engines, or reporting vintages.

The software interface consequently needs an explicit requested horizon,
an observed-support horizon, an event convention and a monetary payment clock.
A response can legitimately contain a point estimate on the supported range
and a set beyond it. It must distinguish extrapolated uncertainty conditional
on a tail model from uncertainty that includes the choice of that model.
Returning a finite number at every requested time does not settle the
identification question.

\section{Why a late hazard can reverse a policy ranking}
\label{exec17:s5-selection}

Consider two equally weighted baseline risk strata, high and low, and no
competing exit or censoring. Within each stratum an active account defaults
with a constant probability per period. Under the reference policy these
probabilities are $0.8$ and $0.1$; under the intervention they are $0.4$ and
$0.05$. One can generate the potential histories using independent period
uniforms within each stratum. Randomized policy assignment, consistency and
the stated event clock then suffice for this illustrative policy comparison.
Its reversal does not require an unmeasured common cause of action and default.

Write $h_b^d$ for the hazard in baseline stratum $b$ and $\pi_b=1/2$ for its
baseline proportion. Marginal survival and the opening-period risk composition
are
\begin{align}
 S^d(j)&=\sum_b\pi_b(1-h_b^d)^j,\notag\\
 w_b^d(j-1)&=\frac{\pi_b(1-h_b^d)^{j-1}}{S^d(j-1)},\notag\\
 h^d(j)&=\sum_bw_b^d(j-1)h_b^d
        =1-\frac{S^d(j)}{S^d(j-1)}.
 \label{exec17:s5-mixture-hazard}
\end{align}
The two marginal hazards average over different survivor populations. By
period three, the reference high-risk share is $4/85$, whereas the policy
high-risk share is $144/505$. Protecting the high-risk accounts keeps more of
them in the later policy risk set. Comparing those risk sets as if they were
the same baseline population confuses composition with a policy's cumulative
consequence.

The exact results make the reversal visible:
\begin{center}
\begin{tabular}{lrr}
\hline
Quantity & Reference & Policy\\
\hline
$S(1)$ & $0.55$ & $0.775$\\
$S(2)$ & $0.425$ & $0.63125$\\
$S(3)$ & $0.3685$ & $0.5366875$\\
Period-three marginal hazard & $113/850$ & $1513/10100$\\
Three-period default risk & $0.6315$ & $0.4633125$\\
Restricted active periods & $1.975$ & $2.40625$\\
\hline
\end{tabular}
\end{center}
The period-three hazards are approximately $0.13294$ and $0.14980$. Thus the
policy has the larger late marginal hazard, although every stratum has a
smaller hazard under that policy at every time and marginal survival is larger
throughout. The marginal hazard ratio is a valid description of its two
selected risk sets. It is insufficient as a stand-alone ranking of the
baseline policy decision.

For a portfolio entering together, suppose each period entered earns USD~20,
default by period three loses USD~1,000, and the policy costs USD~3 upfront.
The incremental expected value is
\begin{equation}
 20(2.40625-1.975)
 -1000(0.4633125-0.6315)-3
 =173.8125\ \mathrm{USD}.
 \label{exec17:s5-selection-value}
\end{equation}
This ranking uses the same initial population and the cumulative event
probabilities appropriate to its financial objective. It remains conditional
on the stipulated margins and losses. Changing those amounts is a sensitivity
analysis of the objective, while changing how histories are selected is a
change to the population or causal model.~\citep{exec17_s5_math}

This example motivates a practical output hierarchy. Preserve stratum-specific
and marginal hazards as diagnostics, but also report the policy survival or
cumulative incidence on a common baseline population, the restricted active
time and their joint uncertainty. If the bank wants an effect among accounts
that would survive under both policies, that is a principal-stratum question
requiring additional assumptions about joint potential histories. Neither
conditioning each fitted curve on its own observed survivors nor adjusting
the plot legend identifies that population.

The same logic affects model assessment. A regression can accurately predict
late defaults among observed reference survivors and still be a poor model
for the histories retained under another policy. Checking calibration on the
observed late risk set is informative about that set. Transporting it to a
policy-induced risk set requires overlap and stability of the relevant
conditional transition law. Flexible prediction and causal transport should
therefore appear as separate arguments in the proposal.

\section{Competing exits and the order of a simulator}
\label{exec17:s5-ordered}

Now distinguish default $D$ and early repayment $P$ as mutually exclusive
first events within each period. Among accounts event-free at the opening,
let $\lambda_{D,j}$ and $\lambda_{P,j}$ be their probabilities conditional on
the permitted history and intervention. They obey
$\lambda_{D,j}+\lambda_{P,j}\leq1$. With constant probabilities in a homogeneous
example,
\begin{align}
 S(j)&=(1-\lambda_D-\lambda_P)^j,\notag\\
 F_D(K)&=\sum_{j=1}^{K}S(j-1)\lambda_D,&
 F_P(K)&=\sum_{j=1}^{K}S(j-1)\lambda_P,\notag\\
 S(K)+F_D(K)+F_P(K)&=1.
 \label{exec17:s5-first-events}
\end{align}
Here $S$ means free of both first events. In contrast,
$1-F_D=S+F_P$ includes accounts that have already repaid. It cannot be
integrated as if those repaid accounts were still earning the active-account
margin. Treating prepayment as censoring would target another quantity and
needs a separate scientific interpretation; it does not automatically
describe eliminating prepayment while keeping every other mechanism fixed.

Some simulators implement within-period draws sequentially. If prepayment is
drawn first with probability $p_P$ and default is then drawn among those
without that prepayment with probability $p_D$, the induced unconditional
first-event probabilities are
\begin{equation}
 \lambda_P=p_P,\qquad \lambda_D=(1-p_P)p_D,
 \qquad P(\text{neither})=(1-p_P)(1-p_D).
 \label{exec17:s5-ordered-conversion}
\end{equation}
Consequently an adapter from mutually exclusive first-event probabilities
must set $p_P=\lambda_P$ and
$p_D=\lambda_D/(1-\lambda_P)$ when $\lambda_P<1$. The constraint on the
two first-event probabilities ensures $p_D\leq1$. If $\lambda_P=1$, feasibility
requires $\lambda_D=0$ and the unused conditional default probability does not
determine any outcome. Reversing the order requires the converse
factorization. Multiplying two arbitrary competing hazards without defining
their conditioning populations supplies neither factorization.

The conversion is elementary but consequential. Consider two periods. Set
$(\lambda_D,\lambda_P)=(0.1,0.1)$ for reference and $(0.12,0.6)$ for policy.
Although the policy's default probability among initially event-free accounts
is larger each period, more of its accounts prepay before a later default is
possible. The resulting $(F_D,F_P,S)$ vectors at period two are respectively
$(0.18,0.18,0.64)$ and $(0.1536,0.768,0.0784)$. The policy has the smaller
default cumulative incidence and the shorter active relationship.

Price each entered period at USD~20, default at a USD~1,000 loss, prepayment
at a USD~50 loss, and the intervention at USD~3 per initial account. The
joint objective is
\begin{equation}
 V^d(K)=20R^d(K)-1000F_D^d(K)-50F_P^d(K)-c_d.
 \label{exec17:s5-event-value}
\end{equation}
Using $c_0=0$ and $c_d=3$ gives reference value $-153$ USD and policy value
$-169.4$ USD. The incremental value is $-16.4$ USD despite a smaller default
incidence. Default avoidance alone does not establish a desirable servicing
policy when relationship duration and competing exits carry monetary value.

For the ordered simulator the correct default probabilities after a
non-prepayment draw are $1/9$ for reference and $0.3$ for policy. Suppose an
adapter instead passes the policy's unconditional $0.12$ directly as $p_D$,
while setting $p_P=0.6$. The induced period default probability is then
$0.4\times0.12=0.048$. At period two the wrong policy calculation returns
\begin{align}
 F_D^{\rm wrong}&=0.064896,&F_P^{\rm wrong}&=0.8112,\notag\\
 R^{\rm wrong}&=1.352,&V^{\rm wrong}&=-81.416\ \mathrm{USD}.
 \label{exec17:s5-ordered-error}
\end{align}
Compared with the correctly translated reference, the apparent gain becomes
$71.584$ USD. An omitted conditional-probability conversion has reversed the
decision. The example does not identify a defect in an upstream package: the
error is the deliberately incompatible meaning supplied at its boundary.
\citep{exec17_s5_math}

The inspected simulation path of pygformula draws a competing event before
the target event and marks the target event missing after that competing
exit. Its aggregation multiplies the target-event probability by the
no-competing-event probability and the preceding survival products. The
stored object named survival is one minus accumulated target incidence;
it is therefore distinct from event-free survival when competing exits are
present. These are claims about the inspected original simulation and
aggregation paths, not about newly fitted models or all package releases.
The model-fitting risk sets were not newly inspected for this comparison,
so the simulation order alone cannot establish how every fitted probability
is conditioned.~\citep{v13_sw_pygformula,exec17_s5_paths}

An integration should first check the complete tuple: opening risk set,
within-period order, event codes, probability conditioning, transition to an
absorbing state, and the meaning of every accumulated output. A generated
two-period law is a useful exact semantic test before statistical fitting.
It can verify that a proposed wrapper conserves probability and obtains
Equation~\eqref{exec17:s5-event-value} under declared inputs. It does not
establish that the observational data identify those inputs, or that a
fitted observational model is calibrated under intervention.

\section{Modifying a natural action and redrawing an action}
\label{exec17:s5-natural}

A dynamic servicing policy may modify the action that would naturally be
chosen: shorten a delay, reduce a proposed payment, or add support when the
proposed contact intensity lies in a specified range. This is more specific
than drawing a new action from a distribution with the same conditional
marginal frequencies. At later times the natural action is evaluated under
the history produced by earlier interventions. Unobserved features of the
action mechanism can connect that later natural action to an earlier one,
even when there is no unmeasured confounding of the ordinary joint treatment
regime and outcome.

The longitudinal modified treatment policy analysis distinguishes a natural
action policy from its stochastic-intervention counterpart. Under its
structural model, standard sequential randomization separates a current
treatment innovation from later covariate and outcome innovations. The
stronger condition also separates it from future treatment innovations.
Theorem~1 identifies the redrawn stochastic intervention under the standard
condition and the natural modified policy under the stronger condition,
together with support. The two targets must therefore be named before
reusing the same observed-data recursion.~\citep{lmtp2023}

An exact two-time construction shows why this distinction is necessary.
Let $U$ and $V$ be independent fair bits, with no measured covariates. In both
worlds $A_1=U$ and the structural outcome is $Y=A_2$. The second natural action
is generated by
\begin{equation}
 \begin{array}{ll}
 \mathcal W_0:& A_2=V,\\
 \mathcal W_1:&
 A_2=\begin{cases}V,&A_1=U,\\1,&A_1\ne U.\end{cases}
 \end{array}
 \label{exec17:s5-two-worlds}
\end{equation}
In ordinary observations $A_1=U$, so the second branch of $\mathcal W_1$ is
never used. In each world all four $(A_1,A_2)$ cells have probability $1/4$,
and $Y=A_2$. Thus the complete observable law is identical, with positive
support for both actions at each observed history. This is equality of the
whole joint distribution, not merely equality of observed outcome means.

Apply the policy $d_1(a)=1-a$, followed by $d_2(a)=a$: flip the first natural
action and leave the second natural action as generated under the changed
history. In $\mathcal W_0$, the second action remains $V$ and
$E(Y^{\bar d})=1/2$. In $\mathcal W_1$, the first action becomes $1-U$;
the second action mechanism then takes its other branch and sets $A_2=1$.
Its modified-policy mean is $1$. The natural-action effect is not identified
by the common observed law under the weaker restrictions of this construction.

For the redrawn intervention, the observed conditional action law at both
times is a fair bit. Pushing those conditional laws through the flip and
identity maps still gives fair bits; drawing from them under the evolving
history yields outcome mean $1/2$ in both worlds. This is a different
intervention. It is identified by the common observed law under the stated
standard condition. Calling it a numerical approximation to the natural
policy would obscure a difference in counterfactual meaning.

The construction satisfies the standard innovation condition used above:
covariate and outcome structural innovations can be constants, and the
outcome under a joint assigned sequence is simply the assigned second
action. The second treatment innovation in $\mathcal W_1$, however, contains
the earlier $U$. Strong independence between the first and later treatment
innovations fails. In $\mathcal W_0$ independent $U,V$ suffice. Support also
holds in the example; the failure does not arise from a missing observed
action cell. These details explain why the witness illustrates the theorem's
assumption distinction rather than contradicting its identification result.
\citep{lmtp2023,exec17_s5_math}

In banking terms, a latent servicing preference that influences several
successive proposed actions can matter to a policy that modifies those
proposals. Observing all outcome confounders for ordinary assigned regimes
does not automatically reveal that proposal mechanism. This interpretation
is a possible mechanism, not an assertion about a bank's data. A proposal
should ask whether the natural recommendation can be recorded, whether its
history-dependent generation can be justified, and whether a randomized
redraw is operationally acceptable. If the business instead intends a fresh
randomized allocation at every time, the stochastic target may be exactly
the intended one.

There are consequently several defensible responses to uncertain natural
action dependence: record and experimentally vary the recommendation
mechanism; impose and defend the stronger structural restriction; redesign
the intervention as a redrawn policy; or retain an explicitly incomplete
identification claim. An estimator's robustness properties concern an
identified observed-data functional under its assumptions. They cannot
choose among observationally identical structural worlds with different
natural-policy outcomes.

\section{A supported shift, a cap, and the boundary of the density ratio}
\label{exec17:s5-pushforward}

Continuous treatment policies create a second boundary that is easy to hide in
an API. Let $A\sim\operatorname{Uniform}[0,1]$ and consider the fixed map
\begin{equation}
 d(a)=\begin{cases}a+1/4,&0\leq a\leq3/4,\\a,&3/4<a\leq1.\end{cases}
 \label{exec17:s5-shift}
\end{equation}
The map is not globally one-to-one. Its image has two inverse branches on
$[3/4,1]$: the identity branch and the shifted branch. The pushforward density
relative to Lebesgue measure is therefore
\begin{equation}
 g_d(y)=\begin{cases}
 0,&0\leq y<1/4,\\
 1,&1/4\leq y<3/4,\\
 2,&3/4\leq y\leq1.
 \end{cases}
 \label{exec17:s5-push-density}
\end{equation}
The value two is not a probability exceeding one. It is the sum of the masses
from two preimages. Dropping one inverse branch changes the intervention law
and its likelihood ratio. Integrating the density gives one, and direct
integration gives
\begin{equation}
 E\{d(A)\}=11/16,\qquad E\{d(A)^2\}=25/48,\qquad
 \operatorname{Var}\{d(A)\}=37/768.
 \label{exec17:s5-push-moments}
\end{equation}

Now compare the cap $d_c(a)=\min(a+1/4,1)$. It maps the
interval $(3/4,1]$ to the single point one, creating an atom of mass $1/4$.
The observed continuous treatment law gives that point probability zero, so
the intervention law is not absolutely continuous with respect to the observed
law. A density-ratio estimator defined only with respect to the observed
Lebesgue density cannot represent this target by an ordinary Radon--Nikodym
ratio. The endpoint belongs to the topological support $[0,1]$; the problem
is the stronger absolute-continuity requirement. A large finite ratio and a
singular atom are different mathematical objects.

The distinction matters even if an outcome regression is available. Suppose
$m(a)=E(Y\mid A=a)$ is known only as a version up to Lebesgue-null sets and
the outcome is bounded in $[0,1]$. The value at the cap point is then
unrestricted by the observed regression. Suppose consistency and
exchangeability identify the causal response mean almost everywhere on that
continuum, without imposing endpoint continuity. The target mean then lies in
\begin{equation}
 \left[\int_0^{3/4}m(a+1/4)\,da,
       \int_0^{3/4}m(a+1/4)\,da+1/4\right]
 = [15/32,23/32]
 \label{exec17:s5-cap-set}
 \end{equation}
for the exact authored regression $m(a)=a$ almost everywhere on the observed
continuum. The interval is sharp. Let $E$ be uniform on $[0,1]$, independent
of $A$, and define $Y^a=1\{E\leq m_u(a)\}$, where $m_u(a)=a$ for $a<1$ and
$m_u(1)=u$ with $u\in[0,1]$. Every such world has the same complete
observed $(A,Y)$ law, because $P(A=1)=0$. Its capped mean is
$15/32+u/4$. The two endpoint worlds and their intermediate values therefore
attain the entire set while preserving exchangeability and consistency.

If the outcome is worth USD~100 at its unit value, the baseline is USD~50 and the
intervention costs USD~3, the net value set is
$[-49/8,151/8]$ USD. A continuous extension of $m$ to the endpoint would
force $u=1$ in this example. A structural model for the atom supplies a
different explicit restriction. Redesigning the intervention to avoid the
atom changes the policy instead. These remedies address distinct scientific
questions; they must not be interchanged by an automatic clipping rule.
\citep{lmtp2023,exec17_s5_math}

Identification under continuity is also distinct from regular estimation.
A value defined at a single continuous treatment point can depend on local
behavior that the nonparametric observed-data model does not estimate at
root-$n$ rate. The LMTP paper's threshold example describes this
non-pathwise-differentiability and restricts its efficiency theorem to
piecewise smoothly invertible maps. Continuity can resolve the choice of an
endpoint version without proving the square-integrable influence function
required by that theorem. A parametric model, quantitative smoothness
restriction, mixed treatment law with a genuine endpoint atom, or a smoothed
policy would need its own estimation argument. The cap is operationally
meaningful; the unrestricted continuous-density theorem is simply not its
automatic uncertainty guarantee.~\citep{lmtp2023}

\section{The derivative of a fixed modified-action map}
\label{exec17:s5-derivative}

For the supported fixed map in Equation~\eqref{exec17:s5-shift}, write
$g(a\mid W)$ for the observed conditional treatment density,
$r(a,W)=g_d(a\mid W)/g(a\mid W)$, and $m(a,W)$ for the conditional outcome
regression. Here $W$ denotes baseline information available before the action.
Assume the relevant change of variables is valid, the map is fixed, and the
score below has finite second moment. In a one-time setting, consistency,
exchangeability given $W$, and the required absolute continuity identify the
following causal mean from the observed law. The one-time target is
\begin{equation}
 \theta(P)=E_W\int m(d(a,W),W)g(a\mid W)\,da.
 \label{exec17:s5-lmtp-target}
\end{equation}
One can derive its gradient directly. For a smooth dominated submodel write
the score as $s=s_W+s_A+s_Y$, with
$s_W=E(s\mid W)$,
$s_A=E(s\mid A,W)-E(s\mid W)$, and
$s_Y=s-E(s\mid A,W)$. The three scores have mean zero in their respective
conditioning populations. Differentiation of the outer baseline law, the
assignment law inside the integral, and the conditional outcome regression
gives
\begin{align}
 \dot\theta={}&E\{(M(W)-\theta)s_W\}\notag\\
 &+E\{(m(d(A,W),W)-M(W))s_A\}\notag\\
 &+E_W\int g_d(a\mid W)
       E\{(Y-m(a,W))s_Y\mid A=a,W\}\,da ,
 \label{exec17:s5-path-parts}
\end{align}
where $M(W)=E\{m(d(A,W),W)\mid W\}$. The last line uses the pushforward
change of variables, including every inverse branch, to move the derivative
of $m(d(a,W),W)$ onto the observed treatment scale. Replacing $g_d$ by $rg$
makes that line the expectation of $r(A,W)(Y-m(A,W))s_Y$.
The fixed-map path derivative thus has the mean-zero influence function
\begin{equation}
 D(P)(O)=r(A,W)\{Y-m(A,W)\}+m(d(A,W),W)-\theta(P).
 \label{exec17:s5-lmtp-if}
\end{equation}
Each summand recovers its corresponding score covariance in
Equation~\eqref{exec17:s5-path-parts}; all cross-covariances vanish by
conditional centering. This establishes the derivative in the stated
observational model. The primary LMTP theorem supplies the longitudinal
generalization under its additional conditions.~\citep{lmtp2023}

The map is fixed independently of $P$, although its induced density $g_d$
changes when $g$ changes. To isolate the distinction, define a second
functional using an external intervention density $q_0(a\mid W)$ held fixed
along the submodel. At the reference law take $q_0=g_d$, so the two targets
have the same numerical value at that law. Under perturbation they need not.
For the external-$q_0$ functional the middle line of
Equation~\eqref{exec17:s5-path-parts} is absent and
\begin{align}
 M(W)&=\int m(a,W)q_0(a\mid W)\,da,\notag\\
 D_{\mathrm{frozen-q}}&=r(A,W)\{Y-m(A,W)\}+M(W)-\theta.
 \label{exec17:s5-frozen-if}
\end{align}
The observed assignment law may still vary; it is the intervention law that
is externally frozen in this comparison. A known randomized observational
assignment defines a restricted statistical model and a different
projection argument. Neither interpretation should be chosen merely from
the fact that a software routine accepts supplied probabilities.

The extra random term in the observational fixed-map score is
$m(d(A,W),W)-M(W)$. Its covariance with the outcome-residual component
vanishes after conditioning on $(A,W)$, and its covariance with the
baseline component vanishes after conditioning on $W$. Consequently
\begin{equation}
 \operatorname{Var}(D)-\operatorname{Var}(D_{\mathrm{frozen-q}})
 =E\!\left[\operatorname{Var}\{m(d(A,W),W)\mid W\}\right]\geq0.
 \label{exec17:s5-assignment-variance}
\end{equation}
This is a comparison of the two specified influence functions at the reference
law, not a claim that their estimators answer the same question away from it.

The exact check uses the shift in Equation~\eqref{exec17:s5-shift} and
$m(a)=a$, with no baseline variation. Along the normalized assignment path
$g_\epsilon(a)=1+\epsilon(2a-1)$ for $|\epsilon|<1$, direct integration gives
\begin{align}
 \theta(\epsilon)
 &=\int_0^{3/4}(a+1/4)\{1+\epsilon(2a-1)\}\,da\notag\\
 &\quad+\int_{3/4}^1 a\{1+\epsilon(2a-1)\}\,da
   =\frac{11}{16}+\frac{23}{192}\epsilon .
 \label{exec17:s5-assignment-path}
\end{align}
Thus $\dot\theta(0)=23/192$ and $\theta(1/4)=551/768$. The covariance
$E[\{d(A)-11/16\}(2A-1)]$ reproduces the derivative; an outcome-residual
score alone gives zero and misses it. Under the external frozen intervention
law, with $m$ held fixed, zero is instead the correct derivative. The check
separates the scientific targets before comparing estimated standard errors.

For a Bernoulli outcome with mean $A$, the residual score variance is
\[
 \int_{1/4}^{3/4}a(1-a)\,da+
 4\int_{3/4}^1a(1-a)\,da=\frac7{32}.
\]
The factor four is the square of the density ratio two, not a second
contribution to the probability mass. Adding the assignment variance
$37/768$ from Equation~\eqref{exec17:s5-push-moments} gives full score
variance $205/768$. For $n$ independent observations an asymptotically
linear estimator with this score has leading variance $205/(768n)$,
subject to its nuisance-rate and remainder conditions. The arithmetic
checks those components but does not establish those conditions for a
fitted learner or an actual bank. The monetary outcome scale multiplies
the score linearly and its variance quadratically.~\citep{exec17_s5_math}

\section{Translating software semantics before comparing estimates}
\label{exec17:s5-software}

The common package names in a longitudinal workflow conceal different data
contracts. The inspected ltmle documentation uses wide, time-ordered nodes;
an absorbing event is encoded as event one, censoring is a separate factor,
and post-event or post-censor nodes are ignored. Its documented bounds act on
cumulative treatment probabilities. A continuous outcome range can also be
truncated or scaled. These choices determine the observed target before a
targeting step is run.~\citep{sw_ltmle}

The lmtp documentation uses chronological treatment, covariate and outcome
nodes. A supplied \texttt{shifted} data frame overrides the \texttt{shift}
function; it is not an additional rule applied after that function. The
finite-memory parameter specifies how many previous time points are used
for estimation. That documented modeling choice does not establish that a
bank's history can be truncated without changing its conditional law.
Cluster identifiers, folds and sampling weights affect the fitted calculation.
Its returned density ratios are noncumulative. Multiplying them to form the
cumulative ratios required by a stated score is valid. Treating each
noncumulative ratio as an already cumulative quantity, or multiplying an
already cumulative output again, changes that score.~\citep{sw_lmtp}

For a two-step exact probability law with incremental denominators
$(1/2,2/5)$, the incremental ratios are $(2,5/2)$ and their cumulative
products are $(2,5)$. If a floor of $1/4$ is applied to the cumulative
denominator, the final weight is $4$; if the floor is applied to each
incremental denominator before multiplying, the final weight is $5$. Both
procedures avoid an infinite value, but they do not implement the same
numerical calculation. The floor level and location belong in the estimator
specification. A fixed floor can leave bias for the original target, while
a vanishing floor needs rate and moment arguments. Clipping a denominator
does not automatically define a different normalized intervention law.
If a deliberately changed target is intended, that target needs its own
definition and identification argument.~\citep{exec17_s5_math}

The stremr manual presents a long-data interface with subject and time fields
and separates treatment, monitoring and censoring models. Its `fitTMLE` path
wraps sequential g-computation with targeting. The survtmle manual fixes a
binary treatment, integer event times and an event code for censoring; its
default horizon is derived from the observed event time. It documents mean and
hazard methods, covariance and subject-level influence vectors, and warns that
finite-sample convergence is not guaranteed. The prose and argument material
also use different descriptions of a tolerance scaling: the displayed call
uses $1/\sqrt n$, while the descriptive paragraph names $1/n$.
An adapter must inspect the selected executable path and record the actual
argument, rather than resolve that discrepancy by choosing whichever
default sounds preferable. Reaching the maximum iteration count and meeting
the empirical-score tolerance are different termination events.
\citep{v12_sw_stremr,v12_sw_survtmle}

The survtmle horizon illustrates why defaults need scientific review.
The saved API uses the maximum failure time among noncensored event records,
not the end of all recorded follow-up. Two samples with identical requested
six-period objectives can therefore have different implicit output times.
Passing an explicit common horizon removes that software difference, but
does not create positivity of observation or identify an unsupported tail.
Likewise, its covariance and subject influence vectors are useful for
post-hoc contrasts only if the requested treatment, event, horizon and
subject indexing remain attached to their entries. The interface itself
does not establish successful targeting or valid nuisance rates.

The gfoRmula API exposes risk and end-of-follow-up outcomes, censoring and
competing-event options, and custom intervention and history roles. Its names
do not by themselves establish whether an output is event-free survival,
cumulative incidence, or a different complement. The pygformula simulation
path draws competing events before target events, while its bootstrap resamples
whole histories for unique subjects, refits models, and uses common seeds for
intervention simulations. Those are useful building blocks for a prospective
comparison. They do not prove that every model-fitting risk set or every
release has the same semantics.~\citep{sw_gformula,v13_sw_pygformula,v131_sw_pygformula}

These packages also answer different design questions. A fixed binary
baseline-treatment survival analysis can be compared with survtmle's
declared input. Repeated binary decisions with censoring and an ordered
wide history fit the selected ltmle contract; a long history with a distinct
monitoring process motivates examining stremr's corresponding treatment,
monitoring and censoring components. Modified continuous action laws motivate
the lmtp map, support and ratio checks developed above. Parametric
g-computation motivates gfoRmula or pygformula when the goal includes
simulating full intervention histories. An adapter can make their results
comparable only on an intersection of targets and observation assumptions.
It cannot turn a fixed baseline treatment into an evolving decision merely
by repeating its value down a long table.

The distinction between ordinary and sequential double robustness also
belongs in this comparison. The inspected sequential-robustness paper
distinguishes consistency when either every outcome regression is correct
or every treatment mechanism is consistently estimated from consistency
when, at each time, the appropriate outcome or treatment condition holds.
Its outcome conditions concern functional specification of the recursively
defined regressions, and its formal rate and inference conditions add
further restrictions. A product should therefore identify which algorithm,
target and recursion support a robustness statement. The name TMLE, an
ensemble learner, or an empirical score near zero cannot by itself establish
the stronger per-time property. No rederivation of the earlier sequential
remainder is needed to make this software choice explicit.
\citep{sequentialdr2018}

An interoperable result should therefore carry at least: subject unit and
cluster, clock and horizon, event and censoring codes, absorbing-state rule,
policy law and whether it is natural or redrawn, treatment probability scale,
ratio orientation and truncation, target type, influence-vector convention,
failed-replicate handling, software version and license. A result lacking one
of these fields can remain a useful exploratory number, but it cannot enter a
common-target comparison without an explicit refusal or translation proof.
Package declarations currently record gfoRmula 1.1.2 (GPL-3), ltmle 1.3-0
(GPL-2) and lmtp 1.5.4 (AGPL-3.0). The inspected API bytes and package
declarations are separate evidence; metadata alone does not prove that those
bytes belong to those exact releases. The selected pygformula source is
pinned to commit \texttt{e66f787fcef2} dated 22 May 2026; the selected
stremr and survtmle records identify 0.4 at \texttt{14d15651b247} and 1.1.1
at \texttt{295de8b4be18}, respectively. These are scoped historical identities,
not a claim that all current releases were surveyed. Full dependency rights
and runtime qualification require separate inspection and trials.
\citep{exec17_s5_paths}

\section{Joint event uncertainty and the monetary objective}
\label{exec17:s5-joint}

The financial objective in Equation~\eqref{exec17:s5-event-value} depends on
three components of the same account history. They should not be treated as
independent measurements. Return to its two-period policy law, with constant
first-event probabilities $d=3/25$ for default, $p=3/5$ for prepayment, and
$s=1-d-p=7/25$ for neither event. Write $B$ for the indicator of no first-period
event, $L=1+B$ for the number of periods entered, and $I_D,I_P$ for indicators
of default and prepayment by period two. The complete subject vector is
$X=(L,I_D,I_P)^\top$. Its mean is
\begin{equation}
 E(X)=\left(1+s,\ d(1+s),\ p(1+s)\right)^\top
      =\left(\frac{32}{25},\frac{96}{625},\frac{96}{125}\right)^\top.
 \label{exec17:s5-joint-mean}
\end{equation}
This follows by conditioning on the first event: an account enters period two
exactly when $B=1$, and can default by the horizon either immediately or after
that event-free first period. No independence between the entries of $X$ has
been assumed.

The covariance has a direct indicator derivation. Since default and prepayment
are mutually exclusive first-event categories, $I_DI_P=0$ account by account.
Thus, writing $f_D=d(1+s)$ and $f_P=p(1+s)$,
\begin{align}
 \operatorname{Var}(L)&=s(1-s),\notag\\
 \operatorname{Cov}(I_D,I_P)&=-f_Df_P,\notag\\
 \operatorname{Cov}(L,I_D)
 &=E(BI_D)-s f_D=sd-sd(1+s)=-ds^2,\notag\\
 \operatorname{Cov}(L,I_P)&=-ps^2.
 \label{exec17:s5-joint-identities}
\end{align}
The remaining diagonal entries are $f_D(1-f_D)$ and $f_P(1-f_P)$.
Whole-history enumeration and these separate identities give the same matrix:
\begin{equation}
 \Sigma=\operatorname{Cov}(X)=
 \begin{pmatrix}
 126/625&-147/15625&-147/3125\\
 -147/15625&50784/390625&-9216/78125\\
 -147/3125&-9216/78125&2784/15625
 \end{pmatrix}.
 \label{exec17:s5-joint-matrix}
\end{equation}
Negative event-indicator covariance expresses exclusion of the alternative
exit. Negative duration--event covariance expresses earlier termination of
the active relationship. Both arise from the same joint history law.

For monetary coefficients $c=(20,-1000,-50)^\top$, the individual net reward
is $Z=c^\top X-3$. The cost is constant, so it affects the mean but not the
variance. The exact moments are
\begin{align}
 E(Z)&=-169.4\ \mathrm{USD},\notag\\
 \operatorname{Var}(Z)&=c^\top\Sigma c
                 =119207.04\ \mathrm{USD}^2,\notag\\
 \sum_{k=1}^3c_k^2\Sigma_{kk}&=130533.12\ \mathrm{USD}^2.
 \label{exec17:s5-joint-money}
\end{align}
Dropping the off-diagonal terms overstates the variance by
$11326.08\ \mathrm{USD}^2$ in this construction. Independence is not necessarily
an anticonservative approximation; its direction depends on the objective's
signs and the joint law. Changing those signs or changing the event mechanism
can change the direction. An integration should preserve the covariance,
rather than assume that discarding it always gives a safer answer.

For $M=1000$ independent complete trajectories drawn from this fixed policy
law, the sample-average net reward has variance $119.20704\ \mathrm{USD}^2$.
This is an exact sampling or Monte Carlo statement for those trajectories.
It is not the variance of a causal estimator fitted to $1000$ observational
subjects with estimated treatment and observation mechanisms. For that
estimator, a valid vector of subject influence contributions $\phi_i$ would
instead supply the covariance appropriate to its fitted target. Under its
own asymptotic-linear and moment conditions, the monetary influence is
$c^\top\phi_i$; a constant intervention cost contributes no influence, while
an estimated cost can contribute another component. Jointly estimated
policy and reference outputs also require their cross-policy covariance.
The covariance identity remains the same, but the random vector must be the
correct one.~\citep{exec17_s5_joint_error}

This distinction guides a substantive adoption choice. When an engine exposes
only marginal standard errors, a monetary combination may be unidentified
from its reported uncertainty even when every point estimate is usable.
Exposing subject-aligned influence vectors or a justified joint bootstrap
can repair that interface. If neither is available, one may derive a
conservative bound from additional covariance restrictions or return
unavailable joint uncertainty. Summing marginal variances asserts an
independence restriction that competing first events visibly violate here.

\section{Sampling, simulation and model error are different quantities}
\label{exec17:s5-uncertainty}

Let $D_n$ denote the full observed sample of independent subject histories,
including the records used for fitting and the empirical baseline population.
Conditional on it, let $U$ contain the simulator's baseline resampling and
transition draws. For a value estimate using $M$ simulated trajectories,
the law of total variance gives
\begin{equation}
 \operatorname{Var}(\widehat V)=
 \operatorname{Var}_{D_n}\{E(\widehat V\mid D_n)\}+
 E_{D_n}\{\operatorname{Var}(\widehat V\mid D_n)\}.
 \label{exec17:s5-total-variance}
\end{equation}
The first term is variation in the fitted baseline and observed sample; the
second is simulation noise conditional on that fitted object. If the simulator
is conditionally unbiased for a fitted-model functional $V(D_n)$, its first
conditional mean is exactly $V(D_n)$. Model misspecification can leave a
bias $E\{V(D_n)\}-V_{\rm target}$ that is not removed by increasing $M$; it
can also affect the two variance terms. Squared-error risk adds the square
of that bias to the displayed variance. A
reported Monte Carlo standard error that omits the first term is not a
confidence interval for the full causal estimand.

For a paired policy contrast, use one common uniform draw per simulated unit
and set the event indicator to $1\{U<p_d\}$. With $p_0=0.2$ and $p_1=0.3$, the
single-draw difference has mean $0.1$ and variance $0.09$, because it is one
when $U\in(0.2,0.3]$ and zero otherwise. Independent draws have variance
$0.2(0.8)+0.3(0.7)=0.37$. At $M=10{,}000$, the corresponding simulation
variances are $0.000009$ and $0.000037$. Common random numbers can improve a
contrast by coupling comparable paths, but resetting a seed alone does not
guarantee that branching histories remain aligned. The coupling rule must be
declared.

If each independent account contributes two identical rows, treating rows as
iid understates the variance of a subject-level mean by a factor of two. The
bootstrap and influence calculation must resample or retain the independent
subject unit. The pygformula bootstrap's whole-history resampling is therefore
semantically different from a row bootstrap. A future package comparison must
record this unit explicitly.~\citep{v131_sw_pygformula,exec17_s5_math}

Coverage also needs a complete denominator. In an illustrative 200-request
study, 160 requests are available, 152 are covered, so availability is $0.8$,
conditional coverage is $152/160=0.95$, and request coverage is $152/200=0.76$.
Reporting only the conditional fraction would hide the failed or unavailable
requests. Pointwise intervals for several horizons or policies are not
simultaneous bands; retaining the joint subject influence vectors permits a
declared covariance calculation, while simultaneous control needs an explicit
multiplier or resampling rule. These are reporting choices with direct business
meaning when a decision is based on the worst horizon or the best policy.

\section{A bounded transition-error envelope}
\label{exec17:s5-error}

An integration can state a conditional transition-error budget while retaining
its scientific premises. Suppose the reference and approximate processes have
the same initial history law on standard Borel spaces, use the same fixed
policy and reward functions, and differ only in transition kernels
$P_t(\cdot\mid h)$ and $\widetilde P_t(\cdot\mid h)$. Include any randomized
policy draw in the common transition description. Define total variation as
$\sup_E|P_t(E\mid h)-\widetilde P_t(E\mid h)|$, and suppose it is at most
$\varepsilon_t\in[0,1]$ uniformly over histories reachable by either process.
These are restrictions on the compared laws, not conclusions from predictive
fit on the observed histories.

Couple the initial histories identically. While their full histories agree,
use a measurable maximal coupling of the next kernels; its conditional
probability of a first split is at most $\varepsilon_t$. Once histories have
split, regard the split flag as permanent, even if current states later agree.
If $a_t$ denotes the actual probability of a split by step $t$, then
$a_{t+1}\leq a_t+(1-a_t)\varepsilon_t$. The map on the right is increasing
in $a_t$. Induction therefore bounds $a_t$ by
\begin{equation}
 q_0=0,\qquad q_{t+1}=q_t+(1-q_t)\varepsilon_t,
 \qquad q_t=1-\prod_{u=0}^{t-1}(1-\varepsilon_u).
 \label{exec17:s5-coupling-recursion}
\end{equation}
The event of unequal histories is contained in the split event; this also
bounds the total-variation distance between their history laws. A simpler
union bound, useful when an additive budget is easier to report, gives
\begin{equation}
 d_0=0,\qquad d_{t+1}\leq\min\{1,d_t+\varepsilon_t\}.
 \label{exec17:s5-tv-recursion}
\end{equation}
One may choose $d_t=\min\{1,\sum_{u<t}\varepsilon_u\}$ as a conservative
bound, with $q_t\leq d_t$. For a common running reward with range at most
$B_t$ on step-$t$ histories, the absolute
expected running-reward bias is at most
\begin{equation}
 \sum_{t=0}^{K-1}B_t q_t\ \leq\ \sum_{t=0}^{K-1}B_t d_t,
 \label{exec17:s5-running-bound}
\end{equation}
and a common terminal loss with range at most $L$ adds at most $Lq_K$,
or conservatively $Ld_K$. To see this,
condition on the common history event and use equality of coupled rewards
there; on the split event the reward difference is bounded by its range.
Summing those expected absolute differences and applying the triangle
inequality proves the stated envelope. If the reward functions themselves
are approximated, their uniform discrepancies add a separate term; the
displayed transition bound does not already include them.

For $K=3$, $\varepsilon_t=0.02$, $B_t=20$ USD and terminal loss range
$L=1000$ USD, the recursion gives
$d=(0,0.02,0.04,0.06)$, whereas the tighter coupling recursion gives
$q=(0,0.02,0.0396,0.058808)$. The additive running-margin and terminal
bounds are USD~1.20 and USD~60, totaling USD~61.20. The tighter components
are USD~1.192 and USD~58.808, totaling USD~60. The difference is a proved
consequence of counting the opportunity for a first split only while the
histories remain coupled.~\citep{exec17_s5_joint_error}

The tighter envelope is attainable under these restrictions. Let the
reference stay in state zero. Let the approximation, while in state zero,
move to an absorbing state one with probability $\varepsilon_t$; otherwise
it stays zero. At state one define both kernels to stay there. The two
kernels differ in total variation by exactly $\varepsilon_t$ at state zero
and agree at state one. The approximate probability of state one at time
$t$ is $q_t$, whereas the reference probability is zero. Choose the common
running reward $B_t1\{\text{state}=1\}$ and terminal reward
$L1\{\text{state}=1\}$. Their expected difference is exactly
$\sum_tB_tq_t+Lq_K$, attaining all components together. This abstract witness
establishes sharpness of the envelope over the declared kernel/reward class.
It does not imply that the stipulated budgets hold for a fitted bank model.
\citep{exec17_s5_joint_error}

A bank can replace these
uniform budgets with validated conservative budgets only after obtaining
appropriate transition audits. A small average error on frequent histories
does not imply this uniform bound on rare but financially important histories.
Likewise, a terminal value range should include all competing event outcomes,
not only defaults, if the policy objective prices them jointly.

The envelope offers a useful admission test for an implementation. The
adapter must expose the horizon, transition decomposition, reward range,
stated TV or coupling budget, and whether the bound includes model fitting,
simulation, and observation error. More precisely, the displayed envelope
bounds the expected reward under two fixed laws; finite Monte Carlo error
and fitted-law sampling uncertainty need the separate arguments above.
For a policy-reference contrast, bounding each policy's law separately and
adding the two envelopes is valid; a smaller paired envelope requires a
justified common coupling. If any budget is unavailable, the system
can return a conditional point estimate with a declared numerical component,
but it should not label the resulting interval a full causal uncertainty
interval. Refusal is a valid output when the requested guarantee exceeds the
declared budget.

\section{A common-target public prototype and adoption routes}
\label{exec17:s5-adoption}

The practical next step is a public or generated comparison with a common
target, not a claim that six packages are interchangeable. Freeze a finite
baseline cohort law, an integer event clock, explicit default/repayment/censor
codes, two policies, a supported horizon, monetary rewards, and a complete
seed and environment manifest. Admit an engine only after exact two-period and
single-shift references reproduce the probability, pushforward and reward
checks above. A candidate that cannot express an atom, competing first event,
natural-versus-redrawn choice, or subject-level resampling is not silently
coerced into the target; it receives a typed incompatibility record.

For each admitted engine, register the input schema, version and license,
training folds, policy specification, simulation coupling, number of baseline
and Monte Carlo draws, interval or band method, failed-replicate rule and
runtime. Compare the policy contrast, restricted active time, each first-event
incidence, joint covariance, availability, width, numerical error and the
resulting monetary decision. Keep model bias as a separate unknown. A generated
potential-outcome law establishes the truth of that construction only. Public
data can exercise the observation and clock pipeline, but cannot establish a
bank's unmeasured assignment truth or accepted loss threshold.

Adoption then has four defensible routes. Direct adoption is suitable when an
engine's target, event semantics, license and runtime are accepted. A thin
adapter is preferable when the estimator is sound but its schema or result
metadata differ. Porting is justified only when the missing capability is
material, the license permits it, and a reference implementation can be
reproduced. Rejection is appropriate for an unsupported target, unresolvable
license or dependency risk, missing uncertainty meaning, or an output that
cannot preserve the bank's information clock. These are proposal decisions,
not claims that the underlying packages are defective.

The shared library should expose a typed analysis manifest and a refusal
object alongside estimates. The CLI can validate the manifest, run exact
semantic fixtures, and emit a versioned result bundle. An MCP surface can
request validation, explanation, and comparison while preserving the same
manifest and audit record; it should not infer a causal diagram or silently
choose a horizon from an unqualified natural-language request. LLM or agent
assistance can propose mappings and diagnostics, but a human or registered
rule must approve the target, assumptions, source clock and policy. The
longitudinal engine should return an identified set or an incompatibility when
the tail, atom, event order or natural-policy condition is unsupported.

A prospective public/generated pilot can test these translations before an
expensive bank experiment. Its adoption decision should distinguish an exact
semantic match from statistical qualification, complete dependency rights,
accepted operating behavior and economic usefulness. Authentic assignment and
observation records, feasible policy versions, independent review and accepted
bank loss thresholds remain necessary for those later decisions.
\citep{exec17_s5_design}

\section{Transition error along the histories that a policy actually visits}
\label{exec17:s5-occupied}

The uniform envelope above has a useful virtue: it remains valid without
knowing where a transition error occurs. Its price is that a large loss range
can be charged at every possible first split. A servicing comparison can be
more informative when it preserves both the histories visited under the
policy and the financial consequences of the particular transitions that
differ. This requires additional information. It cannot be obtained by
renaming an average predictive error as a uniform transition budget.

Consider a finite horizon of $K$ transitions and a finite state $X_t$ containing
the information needed by a fixed policy. The full finite history may be used
as the state; compressing it into a balance and a delinquency band requires a
separate sufficiency assumption. Let $P_t$ be the reference transition kernel
and $Q_t$ its approximation after applying the same policy. They start from
the same distribution $\mu_0$ and use the same running rewards $r_t(x)$ and
terminal reward $g(x)$. Thus the comparison concerns model error for one
policy. It is not itself the causal effect of switching policies. Put
\begin{equation}
 \begin{split}
 v^Q_K(x)&=g(x),\\
 v^Q_t(x)&=r_t(x)+\sum_y Q_t(y\mid x)v^Q_{t+1}(y),\\
 V(Q)&=\sum_x\mu_0(x)v^Q_0(x),\\
 \mu^P_t(x)&=\Pr_P(X_t=x).
 \end{split}
 \label{exec17:s5-occupied-values}
\end{equation}
The corresponding $P$ recursion defines $v^P$ and $V(P)$. The rewards can
include active-account margin and terminal indicators for mutually exclusive
default and prepayment. The latter are sufficient for undiscounted event
losses. Discounted payments or event-time costs need the corresponding
time/state reward representation; an unchanged terminal indicator cannot
silently account for them.

Define the signed continuation residual
\begin{equation}
 b_t(x)=\sum_y\{P_t(y\mid x)-Q_t(y\mid x)\}v^Q_{t+1}(y).
 \label{exec17:s5-occupied-residual}
\end{equation}
Subtracting the two backward recursions gives
$v^P_t-v^Q_t=P_t(v^P_{t+1}-v^Q_{t+1})+b_t$. Average against $\mu^P_t$.
The first term becomes the average of the next value difference against
$\mu^P_{t+1}$. Repeating this step cancels every intermediate difference;
the terminal difference is zero because $g$ is common. Hence the exact identity
is
\begin{equation}
 V(P)-V(Q)=\sum_{t=0}^{K-1}\sum_x\mu^P_t(x)b_t(x).
 \label{exec17:s5-occupied-identity}
\end{equation}
No small-error expansion is used. The approximation's continuation value is
weighted by occupancy under the reference process. Replacing that occupancy
by the approximation's occupancy in this formula is generally incorrect.
There is a second exact identity with $Q$ occupancy and $P$ continuation
values, obtained by reversing the subtraction. It does not make the unknown
reference continuation values available from an approximate simulator.

For a row $(t,x)$, let $\epsilon_t(x)$ bound its total variation and let
$R_t(x)$ be the range of $v^Q_{t+1}$ over a set containing the support of the
signed row difference. That signed difference has total mass zero. Its
positive and negative masses are equal to its total variation, so subtracting
the smallest continuation value proves
\begin{equation}
 |b_t(x)|\leq \epsilon_t(x)R_t(x),\qquad
 |V(P)-V(Q)|\leq
 \sum_{t,x}\mu^P_t(x)\epsilon_t(x)R_t(x).
 \label{exec17:s5-occupied-envelope}
\end{equation}
Using the full next-state space gives a valid range even when the signed
support is unknown. A smaller range is valid only when the excluded
transitions are known to have zero signed discrepancy. Zero observed counts
or equal rounded probability estimates do not establish that restriction.
The financial range is a continuation-value range, rather than merely the
current period's margin. An error that changes a default probability can
alter a terminal loss even when current margin is unchanged.

An exact three-period law illustrates the distinction. Begin in active state
$A$, with absorbing states $D$ and $P$ for default and prepayment. Both laws
earn USD 20 at each active opening and assign terminal rewards
$g(A)=0$, $g(D)=-1000$ and $g(P)=-50$. Their active transition rows are
\begin{equation}
 \begin{array}{c|ccc}
       & A & D & P\\\hline
 P_t(\cdot\mid A)&0.70&0.10&0.20\\
 Q_t(\cdot\mid A)&0.70&0.11&0.19
 \end{array}
 \qquad t=0,1,2.
 \label{exec17:s5-occupied-laws}
\end{equation}
The approximation moves probability $0.01$ from prepayment to default. It
preserves active survival exactly, so both laws have opening occupancies
$(1,0.7,0.49)$ and restricted active time $2.19$. Their period-three event
vectors $(F_D,F_P,S)$ are $(0.219,0.438,0.343)$ and
$(0.2409,0.4161,0.343)$. Their values are respectively $-197.1$ and
$-217.905$ USD. A correct survival curve alone therefore misses a model-value
error of $20.805$ USD per initial account.

The continuation values make the error traceable:
\begin{center}
\begin{tabular}{c r r r r}
\toprule
$t$ & $\mu^P_t(A)$ & $v^Q_{t+1}(A)$ & $b_t(A)$ & full next-state range\\
\midrule
0 & 1 & $-169.15$ & $9.5$ & $950$\\
1 & $0.7$ & $-99.5$ & $9.5$ & $950$\\
2 & $0.49$ & $0$ & $9.5$ & $1000$\\
\bottomrule
\end{tabular}
\end{center}
At every active row the residual is
$(-0.01)(-1000)+(0.01)(-50)=9.5$ USD. Equation
\eqref{exec17:s5-occupied-identity} gives
$9.5(1+0.7+0.49)=20.805$ USD. The full-range occupied envelope is
$0.01\{950+0.7(950)+0.49(1000)\}=21.05$ USD. If the equality of the
active probabilities is part of the admitted model class, the signed error
can involve only $D$ and $P$, whose continuation range is $950$ USD at every
step. The restricted-support envelope then equals the exact $20.805$ USD.

By comparison, the earlier uniform first-split envelope, using $\epsilon=0.01$,
a running range of USD 20 and terminal range of USD 1000, is
\begin{equation}
 20(0+0.01+0.0199)+1000(0.029701)=30.299\ \mathrm{USD}.
 \label{exec17:s5-occupied-comparison}
\end{equation}
Both bounds are valid in this construction. The tighter answer uses
information about occupancy and the destination of the error that the
uniform envelope did not require. It is not a universally better bound in
every process. Complete event-history enumeration, the backward recursion
and the separate geometric-duration calculation give the same reference
values here. These are exact authored laws, without a fitted external
estimator or observed bank transition audit.~\citep{exec17_s5_occupied}

\section{Which histories can validate the transition budget?}
\label{exec17:s5-transport}

A useful error budget must refer to the histories reached under the requested
policy. A held-out prediction sample often has a different history mix.
Even when its outcomes truthfully measure the desired transition, averaging
its errors under the wrong mix can hide financially important failures.
This distribution question remains after any causal identification argument
has supplied the relevant conditional transition law.

At a fixed time, let $\rho_t(x)$ be a validation input distribution and suppose
$\mu^P_t(x)>0$ implies $\rho_t(x)>0$. For the true signed continuation
residual $b_t$, write $w_t(x)=\mu^P_t(x)/\rho_t(x)$. Cauchy--Schwarz gives
\begin{equation}
 \begin{split}
 \left|\sum_x\mu^P_t(x)b_t(x)\right|
 &=\left|\sum_x\rho_t(x)w_t(x)b_t(x)\right|\\
 &\leq C_t\left\{\sum_x\rho_t(x)b_t(x)^2\right\}^{1/2},\qquad
 C_t^2=\sum_x\frac{\mu^P_t(x)^2}{\rho_t(x)}.
 \end{split}
 \label{exec17:s5-transport-bound}
\end{equation}
Thus a validated upper bound on the squared residual can be transported only
with an appropriate occupancy factor. The inequality concerns the true
conditional residual, not an arbitrary fitted prediction loss. Estimating
that residual or its squared norm with noisy transition observations requires
its own statistical control. A small in-sample squared loss cannot simply be
inserted as the right-hand side of a guaranteed interval.

Consider two history categories with validation probabilities $(0.99,0.01)$,
policy occupancies $(0.1,0.9)$ and residuals $(0,100)$ USD. Validation's mean
absolute residual is only USD 1, while the policy's signed value residual is
USD 90. The validation squared norm is $100$ USD$^2$, and
\begin{equation}
 C_t^2=\frac{0.1^2}{0.99}+\frac{0.9^2}{0.01}
      =\frac{8020}{99},\qquad
 C_t\|b_t\|_{2,\rho}=\sqrt{802000/99}\approx90.0056\ \mathrm{USD}.
 \label{exec17:s5-transport-example}
\end{equation}
The transported bound correctly reflects the nearly complete movement of
decision-relevant occupancy into the rare validation category. It does not
interpret the small unweighted USD 1 as a policy-value guarantee. Such a
residual can be generated by changing a transition to a state whose
continuation reward differs by USD 100; the numbers describe a possible
finite transition law, not merely an algebraic pair of incompatible scores.

If validation never visits the second category, the density ratio there is
undefined. No finite value of $C_t$ repairs that lack of support. With target
occupancy $0.9$ and an otherwise unrestricted continuation range of USD 1000,
the unseen category alone can contribute up to USD 900 in absolute error.
The missing information can be addressed by sampling that category, changing
the policy to avoid it, imposing and defending a transition restriction, or
retaining the resulting broad error set. The interface should expose this
choice. Increasing draws from the first category supplies none of the
missing second-category transition information.~\citep{exec17_s5_occupied}

This support failure should be distinguished from two other problems. First,
an observational transition conditional on a chosen action need not equal
the interventional transition needed by the policy if hidden confounding
remains. Reweighting the history mix cannot resolve that scientific failure.
Second, a measured account state can become available only after the
servicing action. Conditioning on its corrected warehouse value can define
the wrong information history even with complete statistical support. A
validation design must establish transition meaning and the availability
clock before it establishes precision for a residual.

These distinctions suggest concrete uses for public and generated data. A
generated finite process can reveal all transition rows and policy
occupancies, so it can test both the exact identity and deliberate support
failures. A public longitudinal cohort can exercise history construction,
censoring, event order and reweighting diagnostics. If its assignment or
measurement process does not identify the requested causal rows, it can still
test those pipeline operations while leaving causal validation open. Public
data do not become bank transition truth through a change of column names.
Conversely, missing bank truth need not delay a useful semantic prototype
with a fully specified generated process.

\section{A finite validation interval and its longitudinal allocation}
\label{exec17:s5-validation}

The exact identity can support a prospective validation interval under
stronger, explicit sampling conditions. Fix $Q$ using data separate from
the validation sample, so its continuation values are fixed when validation
begins. Index the relevant time/state rows by $j=1,\ldots,m$. Suppose the
reference occupancies $w_j=\mu^P_{t_j}(x_j)$ are known, and obtain $n_j$
independent next-state draws $Y_{jk}$ from the specified reference row $P_j$.
The sample sizes are fixed before observing these draws. Independence
between different rows is not needed for the union bound below, but the
within-row bounded independent-draw premise is needed. Define
\begin{equation}
 \widehat b_j=\frac1{n_j}\sum_{k=1}^{n_j}v^Q_{t_j+1}(Y_{jk})
              -\sum_y Q_j(y)v^Q_{t_j+1}(y),\qquad
 \widehat B=\sum_jw_j\widehat b_j.
 \label{exec17:s5-validation-residual}
\end{equation}
Let $R_j$ bound the range of the sampled continuation value under the
reference row. The bounded-score concentration argument used earlier gives
$\Pr(|\widehat b_j-b_j|>R_j\sqrt{\log(2m/\alpha)/(2n_j)})\leq\alpha/m$.
Apply the union bound over the rows and then the triangle inequality with
nonnegative occupancies. With probability at least $1-\alpha$,
\begin{equation}
 \left|\{V(P)-V(Q)\}-\widehat B\right|\leq
 H(\boldsymbol n):=
 \sqrt{\frac{\log(2m/\alpha)}2}
 \sum_{j=1}^m\frac{w_jR_j}{\sqrt{n_j}}.
 \label{exec17:s5-validation-interval}
\end{equation}
Under these premises $V(Q)+\widehat B\pm H$ is a finite reference-value
interval. The result is conditional on the separately trained model and the
known occupancies. It validates the model's value error from genuine
reference transitions; generating extra draws from $Q$ itself estimates a
different expectation and cannot validate $P-Q$.

Known occupancy is a material premise, not a harmless convenience. In the
three-period example it follows from the scientifically imposed active
probability $0.7$ under both laws. The validation class may allow the
default/prepayment allocation to differ by time, while retaining that known
active probability. One can then validate each active row separately without
assuming that its default fraction is constant across times. When occupancy
is estimated instead, its error multiplies the row residual. A valid analysis
needs a joint confidence region or another justified bound for these weights;
plugging estimated weights into Equation \eqref{exec17:s5-validation-interval}
does not preserve its stated guarantee automatically.

The independent draws can be obtained from a generated reference process,
using distinct trajectories or separate conditional-transition draws.
Observational subject histories generally do not provide arbitrary numbers
of independent transitions at chosen histories. Repeated rows from one
account, adaptive stopping after a favorable residual, estimated causal
weights and random survivor counts require corresponding sampling arguments.
They are not covered by declaring $n_j$ to be the number of available rows.
The public prototype can test the finite result in its declared experiment
before proposing the additional assumptions and costs of authentic
validation.

The interval also shows why equal validation effort at every time can be
inefficient. Put $a_j=w_jR_j$. If one row observation costs $c_j>0$ and the
available variable budget is $C$, consider the continuous planning problem
\begin{equation}
 \min_{n_j>0}\ \sum_j a_jn_j^{-1/2}
 \quad\text{subject to}\quad \sum_jc_jn_j=C.
 \label{exec17:s5-validation-design}
\end{equation}
For positive $a_j$ the objective is strictly convex. Its first-order
conditions are $a_j/(2n_j^{3/2})=\lambda c_j$. Solving them and using the
budget equality gives
\begin{equation}
 S=\sum_j a_j^{2/3}c_j^{1/3},\qquad
 n_j^*=\frac{C(a_j/c_j)^{2/3}}S,\qquad
 H^*=\sqrt{\frac{\log(2m/\alpha)}{2C}}\,S^{3/2}.
 \label{exec17:s5-validation-allocation}
\end{equation}
Convexity establishes the global continuous minimum. The $2/3$ exponent
arises from adding row confidence radii, each proportional to $n_j^{-1/2}$.
It need not equal an allocation that minimizes the variance of a different
estimator. Zero-occupancy or zero-range rows require no observations for this
particular value bound, although other scientific and operating checks may
still need them. Minimum row counts, capacities and integer observations
add constraints; the continuous solution is then a planning reference.

In the example, use three active rows, $\alpha=0.05$,
$w=(1,0.7,0.49)$ and full continuation ranges $(950,950,1000)$ USD. With
equal row costs and budget $30000$, equal allocation gives $10000$ draws at
each time and a simultaneous halfwidth of approximately USD $32.5680$.
The continuous allocation is approximately $(12337.92,9726.89,7935.18)$.
The feasible integer choice $(12338,9727,7935)$ uses the same budget and
has halfwidth USD $32.1752$.

For costs $(1,2,4)$ and budget $60000$ in the stipulated cost units, equal
counts of $8571$ spend $59997$ and give halfwidth USD $35.1783$. The
continuous allocation is approximately $(19905.57,9885.99,5080.61)$.
A feasible rounding and local leftover-allocation rule yields
$(19908,9886,5080)$, spends $60000$, and gives halfwidth USD $31.4016$.
This integer rule is feasible; no claim of exact integer optimality follows
from the continuous proof. The cost units describe transition-validation
effort. They do not establish staff prices, implementation costs or bank
return on investment.~\citep{exec17_s5_occupied}

These are prospective radii from fixed counts, not intervals observed in a
new validation experiment. Even if the estimated center happened to equal
the exact USD $20.805$ discrepancy, a USD $31.4016$ halfwidth would not
certify its sign. A useful design must therefore relate the financial
precision requirement to the sampling cost before commissioning expensive
validation. Better use of event-support restrictions or a justified
homogeneous-row model may reduce that cost, while introducing additional
scientific premises. The calculation makes those tradeoffs visible without
pretending that a large nominal sample guarantees a useful decision.

Finally, the interval addresses one fixed policy, a fixed continuation model
and the specified reference transition experiment. Two policy values require
joint control over their rows and their common or separate reference data.
Choosing a policy or fitting $Q$ using the same validation observations
changes the experiment. Numerical approximation of $v^Q$, reward-model
error, unsupported tails, uncertain natural-action semantics and unmeasured
assignment truth also remain distinct. They can receive justified additional
envelopes or explicit unavailable components; none is already included in
the row-sampling radius. An integration that preserves these distinctions
can explain exactly what a public prototype established and which
information an eventual bank experiment would still need.


\section{Conclusion}
\label{exec17:s5-end}

Longitudinal causal software is comparable only after the decision clock,
event law, intervention meaning, support, uncertainty unit and monetary
objective are aligned. The exact constructions show where a plausible scalar
can reverse a decision: an unsupported tail, survivor composition, an ordered
competing-event conversion, a natural action that is not a redraw, a singular
cap, a cumulative-weight translation, or a simulation error presented as a
confidence interval. A product proposal can close these semantic and design
gaps with public/generated fixtures and bounded derivations now. It cannot
claim that those fixtures establish bank truth, general calibration, or final
adoption.
